\section{Kết luận và hướng phát triển}
\label{sec:conclusion}

\subsection{Tổng kết kết quả}
Nhìn lại quá trình nghiên cứu và triển khai, đồ án đã xây dựng hoàn chỉnh giải pháp
tự động hóa suy diễn logic mệnh đề tuân thủ nguyên tắc không suy đoán trên trò chơi
suy luận trinh thám dạng lưới, kết nối liền mạch giữa nền tảng lý thuyết hình thức
với kỹ nghệ phần mềm và đồ họa tương tác ba chiều. Ở tầng biểu diễn tri thức, nhóm
đã thiết lập hệ thống ánh xạ biến song ánh phân tách rõ không gian biến gốc và biến
phụ, đồng thời hiện thực thành công bộ mã hóa dạng chuẩn hội tự động cho sáu họ lời
khai cốt lõi cùng hai mở rộng quan trọng là đại số vùng giao và mạch chẵn lẻ biến
phụ Tseitin tuyến tính giúp thu gọn số mệnh đề hàng trăm lần trên miền rộng (cụ thể
$565$ lần tại $m = 16$). Toàn bộ hệ thống công thức mã hóa này đã được kiểm chứng
tương đương ngữ nghĩa vét cạn trên hàng chục nghìn cấu hình chân trị độc lập, bảo
đảm tính đúng đắn toán học của tầng biểu diễn.

Tiếp nối tầng biểu diễn, nhóm đã tự cài đặt hoàn chỉnh bộ giải thỏa mãn mệnh đề
DPLL tất định hỗ trợ giả định đầu vào và lan truyền đơn vị đến điểm bất động, làm
bộ giải cốt lõi cho tác tử suy diễn bằng phản chứng. Đặc biệt, việc đề xuất và
chứng minh chặt chẽ cơ chế bộ nhớ đệm nghiệm mẫu giới hạn trong phạm vi từng bước
đã giúp hệ thống cắt giảm tới $18{,}75\%$ số lượt gọi bộ giải trên các màn chơi phi tuyến (trung bình $7{,}20\%$ trên toàn tập), trong khi quy trình mệnh
đề chặn cho phép xác thực tính duy nhất của nghiệm chỉ sau đúng hai lượt giải SAT.
Kết quả thực nghiệm trên tiến trình độc lập qua $14$ màn chơi chuẩn từ $3 \times 3$
đến $5 \times 5$ ghi nhận tỷ lệ giải thành công $100\%$ với thời gian trung vị tối
đa chỉ $42{,}84\text{ ms}$, song hành cùng không gian bàn điều tra ba chiều
\textit{Midnight Archive} mang lại phản hồi đóng dấu trực quan ngay trên từng thẻ
nghi phạm.

\subsection{Hướng phát triển}
Từ những phân tích khách quan về giới hạn hệ thống tại Chương
\ref{sec:experiments}, nhóm định hình bốn hướng nghiên cứu nâng cấp có tính khả thi
cao trong tương lai. Hướng ưu tiên hàng đầu là tái cấu trúc bộ giải mệnh đề theo
kiến trúc học mệnh đề hướng xung đột (\textit{CDCL}) kết hợp cơ chế theo dõi hai
literal (\textit{two-watched literals}) \cite{silva1999grasp,moskewicz2001chaff} và
chiến lược chọn biến dựa trên mức độ tham gia xung đột (như nhánh heuristic VSIDS),
qua đó giảm đáng kể chi phí quét lại các mệnh đề đã thỏa và thao tác sao chép trạng
thái khi quay lui. Song song với cải tiến bộ giải tổng quát, việc bổ sung cơ chế mã
hóa thích nghi cho các ràng buộc đếm số lượng (\textit{cardinality}) --- tự động
chuyển sang mạch bộ đếm tuần tự đa thức $O(m \cdot k)$ của Sinz
\cite{sinz2005towards} mỗi khi kích thước vùng vượt quá tám ô --- sẽ cho phép xử
lý hiệu quả các lời khai bao quát toàn bộ bàn cờ cỡ lớn mà tránh được sự bùng nổ
số mệnh đề.

Bên cạnh hiệu năng trên lưới rộng, để khắc phục hạn chế phân giải của DPLL trước hệ
ràng buộc XOR toàn cục như đã ghi nhận ở kịch bản kiểm thử khắc nghiệt $\text{K04}$,
bộ giải có thể được tích hợp thêm mô-đun đại số tuyến tính thực hiện phép khử
Gauss--Jordan trực tiếp trên ma trận hệ số trường $\mathbb{F}_2$ theo hướng tiếp
cận lai ghép hiện đại \cite{soos2009extending}. Cuối cùng, về mặt khả năng giải
thích trực quan hỗ trợ người học và người chơi, nhóm hướng tới việc khai thác kỹ
thuật trích xuất tập con không thỏa tối tiểu (\textit{Minimal Unsatisfiable
Subset}) từ vết chứng minh vô nghiệm của $KB_t \wedge \neg \ell^*$; khi đó, công cụ
gợi ý trên bàn điều tra ba chiều không những xác định chính xác nghi phạm kế tiếp
có thể chứng minh mà còn tự động tổng hợp được lời giải thích ngôn ngữ tự nhiên
ngắn gọn chỉ dẫn đúng hai hoặc ba lời khai then chốt dẫn tới kết luận ấy.
